\documentclass[12pt]{article}
\usepackage[utf8]{inputenc}
\usepackage[T1]{fontenc}
\usepackage{hyperref}
\title{Zhuang, Zuojiang Language Resources for K-12 Curriculum Development}
\author{Compiled by Claude Sonnet 5.5 (claude-sonnet-5-5) for Andr\'as Kornai}
\date{}
\begin{document}
\maketitle

\section{Introduction}
Zuojiang Zhuang (ISO 639-3 zzj) is a dialect bundle of the Zhuang macrolanguage, spoken along the Zuo River in Guangxi, China, with small communities in Yunnan and in Vietnam. Joshua Project lists about 1.7 million speakers in China. Sources differ on whether it is a single language. The search found no school materials specific to Zuojiang Zhuang. Only reference, audio and community-profile pages turned up.

\section{Primary Educational Resources}
The search found no primers, textbooks or school programmes specific to Zuojiang Zhuang. Standard Zhuang school materials may exist, but the search did not find them.

\section{Linguistic and Reference Materials}
\begin{itemize}
\item Wikipedia article on the Zuojiang Zhuang languages. \url{https://en.wikipedia.org/wiki/Zuojiang_Zhuang_languages}
\item Glottolog entry for zzj. \url{https://glottolog.org/resource/languoid/iso/zzj}
\end{itemize}

\section{Religious and Community Resources}
\begin{itemize}
\item Joshua Project profile of Zuojiang Zhuang in China, which reports Bible portions but no complete Bible or New Testament. \url{https://www.joshuaproject.net/index.php/people_groups/19687/CH}
\end{itemize}

\section{Digital and Audio Resources}
\begin{itemize}
\item Global Recordings Network, Zhuang, Zuojiang. \url{https://globalrecordings.net/en/language/zzj}
\item Global Recordings Network, Ningming Nakan variety. \url{https://globalrecordings.net/en/language/6131}
\item Global Recordings Network, Ningming Mingjiang variety. \url{https://globalrecordings.net/en/language/6133}
\end{itemize}

\section{Curriculum Development Context}
Zuojiang Zhuang has a large speaker base but is not known to be taught in schools. The search found no dedicated educational materials. A curriculum would likely have to draw on the wider Zhuang standard and its writing system, which this search did not cover, and on new community-developed content.

\end{document}